\section{Test functions, approximation and Bochner measures}
\label{sec:preliminaries}

This section records the original-space calculus needed to change
measure. The key points are a compact core for the Laplacian and the
admissibility of Lipschitz multipliers in the Bochner inequality.
All statements concern a full-support \(\RCD(K,N)\) space with
\(1<N<\infty\). They do not involve a weighted curvature bound.

\subsection{A compact core and local approximation}

Besides the class \(\Test(M)\) in \eqref{eq:test-class}, write
\[
\Test^\infty(M)=\{g\in\Test(M):\Delta g\in L^\infty(M)\}.
\]
The superscript denotes a bounded Laplacian. Both test classes are
algebras. Test functions have \(L^2\) Hessians, and the product rule
gives \(\Gamma(g)\in W^{1,2}\) for \(g\in\Test\).
The connection is metric-compatible and torsion-free, with
\(\nabla_{\nabla g}\nabla g=\nabla(\Gamma(g)/2)\)
\cite{Gigli2018,Savare2014}. We denote the \(L^2\) Hessian by
\(H_g\). The graph norm of \(\Delta\) is
\(\|g\|_2+\|\Delta g\|_2\).

\begin{lemma}[Cutoffs and approximation]
\label{pre:approximation}
The following approximations are available.
\begin{enumerate}
\item For a compact set \(E\) in an open set \(U\), there is
\(\chi\in\Test^\infty\) with \(0\le\chi\le1\), equal to one on a
neighborhood of \(E\), and compactly supported in \(U\).
For a fixed basepoint \(o\in M\), there are also exhausting
cutoffs \(\chi_R\in\Test^\infty\),
equal to one on \(B_R(o)\), supported in \(B_{R+1}(o)\), whose
gradient and Laplacian bounds are independent of \(R\).
\item For \(g\in D(\Delta)\),
\begin{equation}
\label{eq:pre-hessian-estimate}
\|H_g\|_2^2\le
\|\Delta g\|_2^2+\max\{-K,0\}\|\nabla g\|_2^2.
\end{equation}
Compactly supported functions in \(\Test^\infty\) form a core
for \(\Delta\).
\item On a fixed compact region, a function \(g\in\Test\) can be
approximated by compactly supported \(g_j\in\Test^\infty\) with
uniformly bounded gradients, with
\[
g_j\to g,\qquad \nabla g_j\to\nabla g,\qquad
\Delta g_j\to\Delta g,\qquad H_{g_j}\to H_g
\quad\hbox{strongly in }L^2
\]
on that region.
\item A compactly supported Lipschitz function \(\xi\) can be
approximated by compactly supported test functions \(\xi_j\)
with a common compact support and a common Lipschitz bound,
uniformly and strongly in \(W^{1,2}\).
If \(\xi\ge0\), the approximants may be chosen nonnegative.
\end{enumerate}
\end{lemma}

\begin{proof}
The good-cutoff construction and the integrated Bochner estimate give
\textup{(i)} and \eqref{eq:pre-hessian-estimate}; see
\cite[Appendix~A, Lemma~A.2 and Equation~(A.13)]{GigliTamanini2021}.
For the local version of the cutoff construction one uses finitely
many small good cutoffs inside \(U\), and a smooth truncation of their
sum. This keeps the Laplacian bounded and Sobolev by the test
algebra and chain rules.

Let \(P_t\) denote the heat semigroup. For \(h\in L^2\cap L^\infty\),
positive-time regularization gives
\(P_th\in\Test^\infty\)
\cite[Appendix~A, Equation~(A.7)]{GigliTamanini2021}.
Spectral calculus and density of bounded \(L^2\) data show that these
regularizations form a core. For each such regularization \(h\),
\[
\Delta(\chi_Rh)
=\chi_R\Delta h+h\Delta\chi_R
+2\langle\nabla\chi_R,\nabla h\rangle.
\]
The uniform cutoff bounds and the \(L^2\) tails imply convergence
to \(\Delta h\) in \(L^2\), as well as convergence of the functions
and gradients. This proves the compact-core assertion.

For \textup{(iii)}, choose nested cutoffs \(\chi,\zeta\) such that
\(\chi=1\) near the region in question and \(\zeta=1\) near
\(\supp\chi\). Put
\[
g_t=\chi P_t(\zeta g).
\]
The function \(\zeta g\) lies in \(\Test\). The product rule,
spectral continuity and the fixed cutoff bounds give convergence
of \(g_t\) to \(\chi\zeta g\) in the Laplacian graph norm.
The Bakry gradient estimate gives
\[
\|\nabla g_t\|_\infty
\le \|\chi\|_\infty e^{-Kt}\|\nabla(\zeta g)\|_\infty
+\|\nabla\chi\|_\infty\|\zeta g\|_\infty.
\]
Apply \eqref{eq:pre-hessian-estimate} to
\(g_t-\chi\zeta g\) to obtain strong Hessian convergence.
Since \(\chi\zeta g=g\) near the region, \textup{(iii)} follows.

For \textup{(iv)}, use \(\xi_t=\chi P_t\xi\), with an outer
cutoff \(\chi=1\) near \(\supp\xi\).
The heat semigroup converges strongly in \(W^{1,2}\);
the uniform Lipschitz estimate and the small-time heat-kernel
moment bound give uniform convergence on the common compact
support. The cutoff product gives the same conclusions for
\(\xi_t\), and positivity is preserved.
\end{proof}


The compact core in Lemma~\ref{pre:approximation} is constructed
using only the original heat flow. It will therefore remain
available when we compare the two generator domains in
Section~\ref{sec:drift}.

\subsection{Weak Hessians and admissible Bochner multipliers}

The weak Hessian in \eqref{eq:weak-hessian} is continuous under weak
local \(W^{1,2}\) convergence of \(F\), with the test gradient and
weight fixed. For \(F\in D(\Delta)\), the integration-by-parts
formula for its \(L^2\) Hessian identifies it with that tensor.
An upper bound for \(H_F(Y,Y)\) means an inequality on every
nonnegative compactly supported Lipschitz weight; it does not
presuppose an \(L^2\) Hessian for \(F\).

The original test class in \eqref{eq:test-class} is contained in
Savar\'e's class \(\mathbb D_\infty\), which requires
\(\Delta g\in W^{1,2}\) but not a bounded Laplacian
\cite[Equation~(3.4)]{Savare2014}.
His measure-valued Bochner calculus
\cite[Lemma~3.2 and Equations~(3.7)--(3.9)]{Savare2014},
together with the Hessian-square estimate
\cite[Theorem~3.3.8]{Gigli2018}, gives
\begin{equation}
\label{eq:bochner-hessian}
\Gamma_2(g):=
\Delta\bigl(\Gamma(g)/2\bigr)
-\Gamma(g,\Delta g)\,m
\ge \bigl(K\Gamma(g)+|H_g|^2\bigr)m.
\end{equation}
The first Laplacian in this expression is measure-valued.
The associated nonnegative Bochner remainder is locally finite
and does not charge sets of zero \(2\)-capacity. Here
\[
\Cap_2(E)=
\inf_{\substack{U\supset E\\U\ \mathrm{open}}}
\ \inf_{\substack{v\in W^{1,2}(M)\\v\ge1\ m\text{-a.e.\ on }U}}
\int_M\bigl(v^2+|\nabla v|^2\bigr)\dd m.
\]
Sobolev multipliers are evaluated using their quasi-continuous
representatives, which are defined outside a set of zero capacity.

In particular, a nonnegative compactly supported Lipschitz
multiplier is admissible in \eqref{eq:bochner-hessian}.
Indeed, approximate it by the nonnegative test functions in
Lemma~\ref{pre:approximation}\textup{(iv)}. Strong
\(W^{1,2}\) convergence gives, along a subsequence, convergence
outside a set of zero capacity. The approximants have a common
bound and compact support, so dominated convergence applies
against the fixed local Bochner measure. In the integration by
parts, strong convergence of the multiplier gradients pairs
with \(\nabla\Gamma(g)\in L^2\). This gives the same identity
and inequality for the Lipschitz multiplier.

\subsection{Metric and heat-flow facts}

A finite-dimensional full-support \(\RCD\) space is proper and
geodesic, has positive finite measure on every nonempty bounded
ball, and has the Sobolev-to-Lipschitz property. In particular, a
Sobolev function with essentially bounded weak gradient has a
Lipschitz representative, with Lipschitz constant bounded by that
gradient bound. Bishop--Gromov volume control implies Gaussian
volume growth: for a fixed basepoint \(o\), there is a constant
\(C>0\) such that
\(m(B_r(o))\le C\exp(Cr^2)\) for \(r\ge1\).
These properties are unchanged under bounded equivalence of the
measure, except that a weighted curvature bound still requires proof.

The original heat semigroup \(P_t\) is Markovian and strongly
continuous on \(W^{1,2}\). Its Bakry gradient and
reverse-Poincar\'e estimates give, for bounded data and positive
times, the bounds used in the proof of
Lemma~\ref{weight:resolvent}. We use the standard estimates in
\cite{Savare2014,AGS2015}; the finite-dimensional heat regularity
and cutoff statements used above are collected in
\cite[Appendix~A]{GigliTamanini2021}.
